\documentclass[11pt]{article}
\usepackage[margin=1in]{geometry}
\usepackage{mathpazo}
\usepackage{amsmath,amssymb}
\usepackage{microtype}
\usepackage{booktabs}
\usepackage{graphicx}
\usepackage[dvipsnames]{xcolor}
\usepackage{tikz}
\usetikzlibrary{positioning,arrows.meta,calc,fit,backgrounds}
\usepackage[most]{tcolorbox}
\usepackage{enumitem}
\usepackage{listings}
\usepackage[font=small,labelfont=bf]{caption}
\usepackage{hyperref}
\hypersetup{colorlinks=true,linkcolor=MidnightBlue,citecolor=MidnightBlue,urlcolor=MidnightBlue}
\usepackage{cleveref}
\urlstyle{tt}
\def\UrlBreaks{\do\_\do/\do.}
\newcommand{\file}[1]{\nolinkurl{#1}}
\newcommand{\code}[1]{\texttt{#1}}
\newcommand{\mhh}{m_{hh}}
\newcommand{\bw}{\boldsymbol{\omega}}
\newcommand{\bp}{\boldsymbol{p}}
\newcommand{\bx}{\boldsymbol{x}}
\newcommand{\data}{\mathcal{D}}

\definecolor{trainable}{HTML}{1B6CA8}
\definecolor{physics}{HTML}{2E933C}
\definecolor{datac}{HTML}{7B4EA3}
\definecolor{warn}{HTML}{C0392B}

\newtcolorbox{keyidea}[1][]{enhanced,breakable,colback=physics!6,colframe=physics!70!black,
  fonttitle=\bfseries,title={#1},boxrule=0.6pt,arc=2pt,left=6pt,right=6pt}
\newtcolorbox{aside}[1][]{enhanced,breakable,colback=trainable!5,colframe=trainable!70!black,
  fonttitle=\bfseries,title={#1},boxrule=0.6pt,arc=2pt,left=6pt,right=6pt}
\newtcolorbox{assume}[1][]{enhanced,breakable,colback=datac!5,colframe=datac!70!black,
  fonttitle=\bfseries,title={#1},boxrule=0.6pt,arc=2pt,left=6pt,right=6pt}
\newtcolorbox{watch}[1][]{enhanced,breakable,colback=warn!5,colframe=warn!70!black,
  fonttitle=\bfseries,title={#1},boxrule=0.6pt,arc=2pt,left=6pt,right=6pt}
\setlist{itemsep=2pt,topsep=3pt}
\setlength{\emergencystretch}{3em}

\lstset{language=Python, basicstyle=\ttfamily\small, columns=fullflexible,
  keepspaces=true, commentstyle=\color{physics!70!black}\itshape,
  keywordstyle=\color{trainable!90!black}\bfseries, stringstyle=\color{datac},
  showstringspaces=false, frame=single, rulecolor=\color{gray!40},
  backgroundcolor=\color{gray!3}, xleftmargin=4pt, xrightmargin=4pt,
  aboveskip=6pt, belowskip=6pt}

\usepackage{array,tabularx}
\newcommand{\bbeta}{\boldsymbol{\beta}}
\newcommand{\bC}{\boldsymbol{C}}

\title{\textbf{Has the smooth $23\times80$ model been done before?}\\[4pt]
\large A review of prior work on the parametrisation of $d\sigma/dm_{hh}$ in HEFT,\\
EFT morphing, machine-learned surrogates, and smooth coefficient functions}
\author{P.~Bhat, N.~De~Filippis, K.\,F.~Johnson, H.\,B.~Prosper, R.~Starnes}
\date{October 31, 2026}

\begin{document}
\maketitle

\begin{abstract}
\noindent
Our primary model writes $d\sigma/dm_{hh} = \sum_{i=1}^{23} C_i(\theta)\,a_i(m_{hh})$, with
each coefficient function $a_i$ a penalised cubic spline times a threshold factor, fitted to
the \emph{bin integrals} of 3,989 POWHEG spectra with the exact Poisson variances. The
smoothness of each $a_i$ is set by the evidence, the posterior is exactly Gaussian, and a
discrepancy term grows with the degree of cancellation. Because the model is so simple, we
asked whether it has been done before. We reviewed the di-Higgs EFT literature, ATLAS and
CMS morphing and reweighting, machine-learned surrogates, and the statistics literature. For
each work we record what was done, how the coefficients were obtained, and what it means for
us. \textbf{Every ingredient has a precedent; we found no instance of the combination.} The
23-term form and per-bin coefficients are the community standard; the per-bin weighted
least-squares fit is, in substance, our free $23\times80$ fit. In statistics, our model is a
\emph{varying-coefficient model} estimated as a \emph{penalised composite link model}. We
found no work that makes the HEFT coefficients smooth functions of $m_{hh}$, fits them to bin
integrals, or attaches a cancellation-dependent model error. We recommend presenting the
model as established statistics applied to a problem where the field still uses per-bin
tables, and offering the smooth $a_i(m)$ as a replacement for those tables.
\end{abstract}

\tableofcontents

% ===========================================================================
\section{Our model, for reference}
\label{sec:ours}

So that each prior work can be compared term by term, we restate the model.

\begin{keyidea}[The smooth $23\times80$ model]
\begin{align}
\frac{d\sigma}{dm}(m;\theta) &= \sum_{i=1}^{23} C_i(\theta)\, a_i(m),
  \qquad a_i(m) = \phi(m)\sum_{j=1}^{40}\beta_{ij}\,B_j(m),
  \qquad \phi(m)=\sqrt{1-4m_h^2/m^2}, \label{eq:ours}\\
\sigma_k(\theta) &= \int_{e_k}^{e_{k+1}}\frac{d\sigma}{dm}\,dm
  = \sum_{i,j} C_i(\theta)\,\beta_{ij}\,I_{jk},
  \qquad I_{jk} = \int_{e_k}^{e_{k+1}}\phi\,B_j\,dm. \label{eq:bins}
\end{align}
\begin{itemize}[leftmargin=12pt]
  \item $C_i(\theta)$: the 23 monomials in $\theta=(\kappa_\lambda,c_t,c_{tt},c_{ggh},
    c_{gghh})$ of the NLO HEFT cross section (\cref{eq:buchalla}). Known.
  \item $B_j$: 40 cubic B-splines on 250--1500~GeV, knots uniform in $\sqrt{m-250}$.
  \item Likelihood: Gaussian, with the Poisson variance $q\,\max(f,q)$ of the model
    prediction $f$, iterated (IRLS) to remove the Neyman-$\chi^2$ bias; $q$ is each spectrum's
    weight per event, recovered from the integer counts.
  \item Prior: second differences of each $\beta_{i\cdot}$ Gaussian with precision
    $\lambda_i$; the 23 $\lambda_i$ maximise the evidence.
  \item Posterior: exactly Gaussian in the 920 $\beta_{ij}$; $\gamma\approx602$ effective
    parameters.
  \item Model error: $\varepsilon(m)\sum_i|C_ia_i|$ added in quadrature, with $\ln\varepsilon$
    quadratic in $m$; pull s.d.\ 1.00 on 1,000 held-out spectra.
\end{itemize}
The \emph{free} $23\times80$ fit is the limit with one box function per bin and no prior:
23 independent numbers per bin.
\end{keyidea}

\noindent
The question for this review is therefore four questions:
\begin{enumerate}[label=(Q\arabic*)]
  \item Has anyone fitted the 23 coefficients \emph{per bin} from many runs? (the free fit)
  \item Has anyone made them \emph{smooth functions} of $m_{hh}$?
  \item Has anyone fitted smooth functions to \emph{bin integrals} of counts, with the
    smoothness chosen from the data?
  \item Has anyone modelled the \emph{error of the polynomial form itself}, and found it
    related to cancellation?
\end{enumerate}

% ===========================================================================
\section{How we searched}

We ran web searches (October 2026) on arXiv, journal sites, CERN and Fermilab Indico, and
the statistics literature. We then read the method sections of the key papers, and the
appendices where the coefficients are defined. The searches covered: the NLO HEFT
$gg\to hh$ coefficients and how they are obtained; HH benchmark and reweighting papers; EFT
morphing in ATLAS and CMS; ``EFT-aware'' histograms; neural surrogates and simulation-based
inference for di-Higgs; splines and Gaussian processes for EFT predictions; spline fits to
histograms; composite link models; and varying-coefficient models.

\begin{watch}[Limits of the search]
A web search is not exhaustive. Experiment-internal notes are not public, and a smooth fit
could be buried in a thesis or in the ancillary files of a paper. Before submission we should
(i) run an INSPIRE ``cited by'' search on refs.~\cite{Buchalla2018,Heinrich2020,Alasfar2023},
(ii) read the HH reweighting talks at the LHC Higgs WG (e.g.\ the 2025 WG4 meeting on
``HH signal reweighting for EFT interpretations''), and (iii) ask the HH subgroup conveners
directly whether anyone has tried a continuous $m_{hh}$ parametrisation.
\end{watch}

% ===========================================================================
\section{Di-Higgs in HEFT: the 23-term form and its coefficients}
\label{sec:hh}

\subsection{Buchalla, Capozi, Celis, Heinrich, Scyboz (2018) \cite{Buchalla2018}}
\label{sec:buchalla}

\paragraph{What they did.} First NLO QCD calculation of $gg\to hh$ in the non-linear EFT
(HEFT) with full top-mass dependence. They showed that the NLO cross section is a polynomial
with 23 terms in the five couplings:
\begin{align}
\sigma^{\rm NLO} = {}&A_1c_t^4 + A_2c_{tt}^2 + A_3c_t^2c_{hhh}^2 + A_4c_{ggh}^2c_{hhh}^2
 + A_5c_{gghh}^2 + A_6c_{tt}c_t^2 + A_7c_t^3c_{hhh} \nonumber\\
 &+ A_8c_{tt}c_tc_{hhh} + A_9c_{tt}c_{ggh}c_{hhh} + A_{10}c_{tt}c_{gghh}
 + A_{11}c_t^2c_{ggh}c_{hhh} + A_{12}c_t^2c_{gghh} \nonumber\\
 &+ A_{13}c_tc_{hhh}^2c_{ggh} + A_{14}c_tc_{hhh}c_{gghh} + A_{15}c_{ggh}c_{hhh}c_{gghh}
 \qquad\text{(the 15 LO terms)} \nonumber\\
 &+ A_{16}c_t^3c_{ggh} + A_{17}c_tc_{tt}c_{ggh} + A_{18}c_tc_{ggh}^2c_{hhh}
 + A_{19}c_tc_{ggh}c_{gghh} + A_{20}c_t^2c_{ggh}^2 \nonumber\\
 &+ A_{21}c_{tt}c_{ggh}^2 + A_{22}c_{ggh}^3c_{hhh} + A_{23}c_{ggh}^2c_{gghh}.
 \label{eq:buchalla}
\end{align}
The first 15 terms are the LO polynomial; the eight terms $A_{16}$--$A_{23}$, all
containing $c_{ggh}$, first appear at NLO.

\paragraph{How the coefficients were obtained.} For the total cross section the $A_i$ were
obtained ``via projections and performing a fit'', and the two agree; only the projections
are tabulated, with errors propagated from the input cross sections (correlations
neglected). For $d\sigma/dm_{hh}$ (their appendix A.1): ``The coefficients have been
determined by evaluating 23 differential cross sections, obtained by changing the values of
the five coupling parameters. Then the system of 23 equations is solved to extract the value
of each coefficient.'' Bins are 20~GeV wide, 40 bins with centres 250--1030~GeV, in fb/GeV,
published as ancillary CSV files.

\paragraph{What it means for us.}
\begin{itemize}
  \item \Cref{eq:buchalla} is our $\sum_i C_i(\theta)a_i$; the $C_i$ are these monomials.
    Our ordering and normalisation should be stated relative to theirs in the paper.
  \item Their differential coefficients come from a \emph{square} $23\times23$ solve per bin:
    exact interpolation through 23 runs, with no redundancy. Any noise in the 23 runs goes
    straight into the $A_i$, amplified by the condition number of the coupling matrix. Our
    fit uses 3,989 runs, so it is heavily over-determined and the noise averages down.
  \item The $m_{hh}$ dependence is piecewise constant in 20~GeV bins; nothing is smooth.
  \item \textbf{(Q1)} partly; \textbf{(Q2)--(Q4)} no.
\end{itemize}

\subsection{Capozi and Heinrich (2020) \cite{Capozi2020}}

\paragraph{What they did.} A shape analysis of the $m_{hh}$ distribution in HEFT: scans of
the coupling space, clustering of shapes, and an autoencoder trained on $10^5$
distributions.

\paragraph{How the coefficients were obtained.} Taken from ref.~\cite{Buchalla2018}:
``evaluated in bins of width 20~GeV from 250~GeV to 1050~GeV, i.e.\ for 40 bins'', provided as
CSV tables in fb/GeV for 13, 14 and 27~TeV. They report: ``The median of the statistical
uncertainties of the differential coefficients'' is at most 3\%, but above $\sim$650~GeV
``some $A_i$ coefficients have uncertainties in the 20--30\% range.''

\paragraph{What it means for us.} This is direct evidence that the per-bin tables are noisy
in the tail. Our smooth $a_i$ share information across bins, and our posterior s.d.\ is
2.9\% of the POWHEG noise. A side-by-side plot of our $a_i(m)$ with their tabulated
$\tilde A_i$ at 13~TeV would make the point well (\cref{sec:actions}).
\textbf{(Q2)--(Q4)} no.

\subsection{Heinrich, Jones, Kerner, Scyboz (2020) \cite{Heinrich2020}}

\paragraph{What they did.} Implemented the NLO HEFT calculation in the POWHEG-BOX:
\code{ggHH} with anomalous couplings. This is the generator of our 3,989 spectra.

\paragraph{How the coefficients enter.} The two-loop virtual amplitude is too slow to
evaluate on the fly, so it is ``computed via a grid which encodes the dependence of the
virtual two-loop amplitude on the kinematic invariants'' $\hat s,\hat t$. Grids are made for
``23 linearly independent sets of couplings'' (6,715 points), and the code derives ``for each
phase-space point, the coefficients $a_1,\dots,a_{23}$ by interpolation''. Once the user
chooses the couplings, ``the 23 grids are combined into one'' in the first POWHEG stage. The
Born and real-emission parts are evaluated exactly. They note that these $a_i$ are those of
the Born--virtual interference, not the $A_i$ of the cross section. The grid uncertainty
is ``below 2\% throughout the whole $m_{hh}$ range, except for the first bin'' (6\%, and
12\% for benchmark 5), because the grid points were generated with SM couplings and the low
$m_{hh}$ region is sparsely populated.

\paragraph{What it means for us.}
\begin{itemize}
  \item They do interpolate smoothly, but in $(\hat s,\hat t)$ and for the amplitude, not for
    the $m_{hh}$ spectrum. Different object, different purpose.
  \item \textbf{A possible source of our model error.} Our spectra are generated with a
    \emph{combined} grid, the real emission is computed exactly, and the POWHEG Sudakov
    factor reweights the Born. The 23-term form holds exactly at fixed order, but each of
    these steps can introduce small, coupling-dependent departures from it. Their 6--12\%
    first-bin grid error is consistent with our larger residuals near threshold (deviance per
    bin 1.57 below 555~GeV).
  \item \textbf{(Q4)} indirectly: they quantify the grid's own error, not departures from the
    polynomial.
\end{itemize}

\subsection{LHC Higgs WG note: Alasfar et al.\ (2023) \cite{Alasfar2023}}
\label{sec:alasfar}

\paragraph{What they did.} The community reference for EFT descriptions of HH production,
in HEFT and SMEFT, with benchmark points, translation between the two, and a reweighting
recipe for the experiments.

\paragraph{How the coefficients were obtained.} With POWHEG-BOX-V2 \code{ggHH} at NLO,
13~TeV, PDF4LHC15, $\mu_R=\mu_F=c\,m_{hh}/2$ with $c\in\{\tfrac12,1,2\}$. The inclusive
\textbf{A} and differential $d\textbf{A}$ coefficients are obtained ``using a weighted least
square fit'' to 63 simulations (62 HEFT points and the SM), following ref.~\cite{Buchalla2018}.
The note defines
\begin{equation}
\frac{d\sigma}{dm_{hh}} = \mathrm{Poly}(\boldsymbol{c}, d\textbf{A}\,|\,m_{hh})
  = \boldsymbol{c}^{\mathsf T}\cdot d\textbf{A}, \qquad
w_{\rm HEFT} = \frac{\mathrm{Poly}(\boldsymbol{c}, d\textbf{A}\,|\,m_{hh})}
  {\mathrm{Poly}(\boldsymbol{c}_{\rm SM}, d\textbf{A}\,|\,m_{hh})},
\label{eq:alasfar}
\end{equation}
with bins of 20~GeV from 250 to 1050~GeV, then 1050--1200 and 1200--1400~GeV. Covariance
matrices for $\textbf{A}$ and $d\textbf{A}$ are provided, with error-propagation formulas, and
three coefficient sets for the scale variations.

\paragraph{Caveats they state.}
\begin{itemize}
  \item Events above 1400~GeV take the weight of the last bin: ``not exactly correct, but it
    is the best available estimate.''
  \item ``The total cross section is better predicted by the inclusive \textbf{A}
    coefficients'', so the inclusive set should fix the normalisation.
  \item The $m_{hh}$ shapes close well; $p_T(h)$ shows discrepancies up to $\sim$40\%
    (benchmark 1) and 10--20\% elsewhere, because only $m_{hh}$ is used.
\end{itemize}

\paragraph{What it means for us.}
\begin{itemize}
  \item \textbf{(Q1) yes.} This is our free $23\times80$ fit in substance: per-bin weighted
    least squares over many POWHEG runs, with a per-bin covariance. We must cite it as the
    baseline. Differences: we have 3,989 runs (63 times more), 15~GeV bins to 1455~GeV, the
    exact Poisson variance, and IRLS rather than weights from the observed counts. (The note
    does not define its weights; if they are the observed variances, the fit carries the
    Neyman bias we found and corrected, \cref{sec:baker}.)
  \item Their need for separate inclusive coefficients is a symptom of the binned
    representation. A density model integrates to the total cross section at no extra cost.
  \item The ``last bin'' rule above 1400~GeV and the coarse last two bins are exactly the
    problems a smooth density solves inside its range. (Our own model should still not be
    used beyond 1455~GeV.)
  \item \textbf{(Q2)--(Q4)} no.
\end{itemize}

\subsection{Cadamuro, Ingebretsen Carlson, Sj\"olin (2025) \cite{Cadamuro2025}}

\paragraph{What they did.} Extended the WG reweighting to several kinematic variables at 13
and 13.6~TeV, in HEFT and SMEFT, combining several reference samples with a convex
combination weighted by a distance measure.

\paragraph{How the coefficients were obtained.} ``The coefficients of the polynomials,
\textbf{A}, are derived by performing a fit to MC simulations at parton level'': 63 samples
for HEFT and 70 for SMEFT, couplings drawn uniformly in $[-4\pi,4\pi]$. Again one polynomial
per bin.

\paragraph{Key findings.} Reweighting in $m_{hh}$ alone reproduces $p_T^h$ for constructive
interference but not for destructive interference, with tail discrepancies of $\sim$20\%:
they identify destructive interference as ``the main source for the discrepancies''. They
also note that POWHEG's Sudakov form factor makes the coupling dependence of variables
affected by jet radiation deviate from a strict polynomial.

\paragraph{What it means for us.} This is the closest thing we found to \textbf{(Q4)}. Their
observable and question differ (a second variable, not the $m_{hh}$ fit itself), but the
physics is the same: where the 23 terms cancel, small departures from the polynomial become
large relative errors. Our discrepancy $\varepsilon(m)\sum_i|C_ia_i|$ is a quantitative
model of exactly that, and their Sudakov remark offers a mechanism. We should cite it when
introducing the discrepancy.

\subsection{Carvalho et al.\ (2016, 2017) \cite{Carvalho2016,Carvalho2016b,Carvalho2017}}

\paragraph{What they did.} At LO (15 terms): cluster analysis of $m_{hh}$ and
$|\cos\theta^*|$ shapes to choose the 12 benchmark points \cite{Carvalho2016}; an analytical
parametrisation of the total cross section \cite{Carvalho2016b}; and a per-bin
parametrisation for reinterpreting searches \cite{Carvalho2017}.

\paragraph{How the coefficients were obtained.}
\begin{itemize}
  \item Total cross section \cite{Carvalho2016b}: ``The components of \textbf{A} are
    extracted by maximizing the likelihood simultaneously for all the coefficients'', with
    MINUIT, from 169--266 MadGraph points per energy (10,000 events per point, different
    seeds). The per-point error is set \emph{a posteriori} to 0.03\%, chosen so that the
    residuals have RMS $\approx 1$.
  \item Differential \cite{Carvalho2017}: 59 bins in $m_{hh}$ (10~GeV up to 700~GeV, coarser
    above) by 4 in $|\cos\theta^*|$, one coefficient set per bin, from a simultaneous fit with
    MC statistical errors only (50,000 events per point, 13 million for the SM). They
    note that with only 15 samples per bin the result ``would have a limited reliability'',
    and use lower-dimensional sub-spaces to keep enough events per bin.
\end{itemize}

\paragraph{What it means for us.} The over-determined per-bin fit, \textbf{(Q1)}, goes back
to 2017 at LO. Setting the error a posteriori so that residuals have unit RMS is a crude
form of what our discrepancy term does properly. \textbf{(Q2), (Q3)} no.

\begin{table}[h]
\centering\small
\begin{tabularx}{\linewidth}{@{}l l >{\raggedright\arraybackslash}X
  >{\raggedright\arraybackslash}X@{}}
\toprule
Work & Order & $m_{hh}$ dependence & How coefficients are obtained \\
\midrule
Carvalho 2017 \cite{Carvalho2017} & LO, 15 & 59 $m_{hh}$ $\times$ 4 $\cos\theta^*$ bins
  & simultaneous fit per bin; 50k events/point \\
Buchalla 2018 \cite{Buchalla2018} & NLO, 23 & 40 bins of 20~GeV
  & square $23\times23$ solve per bin from 23 runs \\
Capozi 2020 \cite{Capozi2020} & NLO, 23 & as above & from \cite{Buchalla2018}; 20--30\%
  errors above 650~GeV \\
Heinrich 2020 \cite{Heinrich2020} & NLO & per phase-space point
  & 23 amplitude grids in $(\hat s,\hat t)$, combined linearly \\
Alasfar 2023 \cite{Alasfar2023} & NLO, 23 & 20~GeV to 1050, then 2 wide bins
  & weighted least squares, 63 runs; covariances \\
Cadamuro 2025 \cite{Cadamuro2025} & NLO, 23 & per bin, several variables
  & fit to 63 (HEFT) / 70 (SMEFT) runs \\
\midrule
\textbf{This work} & NLO, 23 & \textbf{smooth} $a_i(m)$, 250--1455~GeV
  & penalised Poisson fit to bin integrals of 3,989 runs; evidence; exact posterior \\
\bottomrule
\end{tabularx}
\caption{How the coefficient functions are represented in the di-Higgs EFT literature.
Before this work the $m_{hh}$ dependence is always piecewise constant.}
\label{tab:hh}
\end{table}

% ===========================================================================
\section{Morphing and reweighting in the experiments}
\label{sec:morph}

\subsection{ATLAS morphing \cite{ATLmorph} and \code{RooLagrangianMorphFunc} \cite{RooMorph}}

\paragraph{What it is.} Any distribution at couplings $\theta$ is written as a linear
combination of $N$ base distributions,
\begin{equation}
T(\theta) = \sum_{n=1}^{N} w_n(\theta)\,T(\theta_n),
\end{equation}
where $N$ equals the number of terms in the polynomial and the weights $w_n(\theta)$ follow
from inverting the square matrix of monomials at the base points $\theta_n$. It is
implemented in ROOT as \code{RooLagrangianMorphFunc}.

\paragraph{What it means for us.} Same algebra as \cref{eq:ours}: the $a_i$ are linear
combinations of base histograms. Morphing applies it to whole histograms, with exactly as many
samples as unknowns, so it is an interpolation, not a fit. There is no smoothing in the
observable, and the statistical errors of distant points are amplified by large weights.
\textbf{(Q1)} no (square system); \textbf{(Q2)--(Q4)} no.

\subsection{Overdetermined morphing: Belyaev, Konoplich, Prokofiev (2025) \cite{Belyaev2025}}

\paragraph{What they did.} Use more base samples than unknowns and solve by least squares via
a QR factorisation, $x = R^{-1}Q^{\mathsf T}y$, to stabilise morphing far from the base
points. They note that standard morphing computes its coefficients without accounting for
the uncertainties of the simulated cross sections; statistical power is tracked by the
effective-weight and Kish ratios.

\paragraph{What it means for us.} This is the morphing community arriving at the
over-determined fit, \textbf{(Q1)}, in 2025. Again no smoothing across bins, and the least
squares is unweighted. Our IRLS with Poisson variances is the statistically correct version
of their step.

\subsection{EFT-aware histograms in CMS \cite{Mohrman2023}}

\paragraph{What it is.} In the CMS top-quark EFT analyses (\code{topcoffea}, \code{histEFT}),
each simulated event carries a quadratic in the Wilson coefficients, fixed by MadGraph
reweighting at enough points and solved with least squares. ``Instead of storing the sum of
weight values, EFT histograms store the sum of the weight parameterizations.''

\paragraph{What it means for us.} An event-level version of the per-bin polynomial; exact
for the generated sample, piecewise constant in the observable. No smoothing; no model
error. Not directly applicable to POWHEG NLO, where per-event reweighting to arbitrary
couplings is not available in the same way.

% ===========================================================================
\section{Machine learning}
\label{sec:ml}

\subsection{Quadratic classifier: Chen, Glioti, Panico, Wulzer (2021) \cite{Chen2021}}

\paragraph{What they did.} For one Wilson coefficient $c$, the ratio of EFT to SM
differential cross sections is a quadratic in $c$ at every point $x$ of phase space. They
build this in:
\begin{equation}
\hat r(x,c) = [1 + c\,n_\alpha(x)]^2 + [c\,n_\beta(x)]^2, \qquad
f(x,c) = \frac{1}{1+\hat r(x,c)},
\end{equation}
with networks $n_\alpha, n_\beta$ trained as a classifier on EFT and SM events at several
values of $c$. ``The quadratic dependence of the differential cross section on the EFT Wilson
coefficients is built-in.'' They contrast it with the parametrised classifier, which takes $c$
as an input and does not enforce the form.

\paragraph{What it means for us.} This is the closest conceptual precedent for HEFTNet and
for any $\sum_i C_i(\theta)a_i(x)$ model: a known polynomial in the couplings multiplying
learned functions of the kinematics. The differences are the target (an unbinned likelihood
ratio from events, versus binned cross sections), the representation (networks versus a
linear spline basis), and the uncertainty (none quoted, versus an exact posterior). Our
finding that, for $m_{hh}$, a 920-parameter linear model matches a network is a useful
counterpoint. \textbf{(Q2)} yes in spirit (smooth functions of $x$ with built-in coupling
dependence), but not for HH, not fitted to bins, and without uncertainties.

\subsection{MadMiner \cite{Brehmer2018} and di-Higgs SBI \cite{Mastandrea2024}}

\paragraph{What they did.} MadMiner uses morphing-aware likelihood-ratio estimators: the
matrix element is ``a polynomial over chosen Wilson coefficients''. Mastandrea, Nachman and
Plehn apply it to di-Higgs in SMEFT with three coefficients, ``a polynomial up to squared
terms, resulting in 10 fit parameters'' and 10 morphing basis points, plus a parametrised
classifier in which $c$ ``is promoted to a feature''. They note reweighting ``works well when
the new parameter point is close to the original one''.

\paragraph{What it means for us.} Unbinned inference, a different goal. Confirms that the
polynomial structure is the standard way to carry coupling dependence; no smooth $m_{hh}$
coefficients.

\subsection{Learned loop amplitudes: Bahl, Braun, Heinrich, Plehn, Revelli (2026) \cite{Bahl2026}}

\paragraph{What they did.} Neural surrogates for the one- and two-loop $gg\to hh$ amplitudes
with full top-mass dependence, with calibrated uncertainties, identification of poorly
described phase-space regions, and adaptive training.

\paragraph{What it means for us.} Complementary: they emulate the amplitude (in
$\hat s,\hat t$) for the SM, not the $m_{hh}$ spectrum across couplings. Their emphasis on
calibrated surrogate uncertainties matches ours. A combination (EFT amplitude surrogates
feeding our spectrum model) is a possible future direction.

\paragraph{Gaussian-process emulators.} GP emulators are common in nuclear EFT (chiral EFT
truncation errors) and cosmology, but we found none for the HH $m_{hh}$ spectrum over the
five HEFT couplings.

% ===========================================================================
\section{Statistics: the model has a name}
\label{sec:stats}

\subsection{Varying-coefficient models: Hastie and Tibshirani (1993) \cite{HastieTibshirani1993}}

\paragraph{What it is.} A regression in which the coefficients of known covariates are
themselves smooth functions of another variable, the ``effect modifier'' $r$:
\begin{equation}
\eta = \beta_0(r) + \sum_i x_i\,\beta_i(r).
\end{equation}
The $\beta_i(r)$ are estimated with smoothing splines or, later, P-splines
\cite{EilersMarx1996,EilersMarx2021}.

\paragraph{The match.} Set $x_i = C_i(\theta)$, $\beta_i = a_i$, $r = m_{hh}$, and
$\eta = d\sigma/dm$: \cref{eq:ours} \emph{is} a varying-coefficient model with 23 covariates,
no intercept, and an identity link. This answers \textbf{(Q2)} at the level of statistical
structure: smooth coefficient functions of a covariate are textbook.

\subsection{Penalised composite link model: Eilers (2007) \cite{Eilers2007}, Rizzi, Gampe,
Eilers (2015) \cite{Rizzi2015}}

\paragraph{What it is.} Counts $y_k$ are observed in coarse groups (age groups, income
brackets). The underlying fine-grained distribution $\gamma$ is smooth. The composite link
model writes the expected counts as
\begin{equation}
\mu = \mathbf{C}\,\gamma, \qquad \gamma = \exp(\mathbf{B}\alpha), \qquad
y_k \sim \mathrm{Poisson}(\mu_k),
\end{equation}
where the composition matrix $\mathbf{C}$ sums the fine grid into the observed groups and
$\mathbf{B}$ is a B-spline basis. The penalised log-likelihood
$\ell(\alpha) - \tfrac{\lambda}{2}\lVert D\alpha\rVert^2$ with a difference penalty $D$ is
maximised by IRLS, and $\lambda$ is chosen by AIC or BIC. Rizzi et al.\ show it recovers
smooth distributions from very coarse, open-ended groups; the R package \code{ungroup}
implements it.

\paragraph{The match.} Our $I_{jk}$ in \cref{eq:bins} is the composition matrix applied to the
spline basis: the model is defined as a density and compared with data through its bin
integrals. Our merged-bin test (fit 30~GeV bins, predict the 15~GeV split) is exactly the
``ungrouping'' that PCLM was designed for. This answers \textbf{(Q3)}. Differences:
\begin{itemize}
  \item \textbf{Link.} PCLM uses a log link so the density is positive. Our $a_i$ change sign,
    so we use the identity link, and the model is linear in $\bbeta$.
  \item \textbf{Covariates.} PCLM smooths one density. We smooth 23 functions, combined with
    known weights $C_i(\theta_s)$ that differ for each of the 3,989 spectra.
  \item \textbf{Smoothing parameters.} AIC/BIC over one $\lambda$, versus 23 $\lambda_i$ by
    maximising the evidence.
  \item \textbf{Uncertainty.} Sandwich or Bayesian intervals for $\gamma$, versus the exact
    posterior of the linear model.
\end{itemize}

\subsection{Older and related ideas}

\begin{itemize}
  \item \textbf{Histosplines} (Boneva, Kendall, Stefanov 1971) \cite{Boneva1971}: splines
    whose integrals over each histogram bin equal the bin content. The first appearance of the
    integral constraint, as interpolation rather than fitting.
  \item \textbf{Evidence and REML for smoothing parameters}: MacKay \cite{MacKay1992} for
    the Bayesian evidence; Wood \cite{Wood2011} for the REML/marginal-likelihood methods that
    are now standard in generalised additive models (\code{mgcv}). Our 23 $\lambda_i$ follow
    this route.
  \item \textbf{P-splines in HEP}: Whitehorn, van Santen, Lafebre \cite{Whitehorn2013} fit
    tensor-product penalised B-splines to IceCube photon tables of up to six dimensions, by
    penalised least squares with one $\lambda$ per dimension (the paper does not choose them
    automatically), with non-negativity and monotonicity where needed. It shows that spline
    surrogates for large Monte Carlo tables are accepted practice in HEP.
\end{itemize}

\subsection{Poisson versus Neyman $\chi^2$: Baker and Cousins (1984) \cite{BakerCousins1984}}
\label{sec:baker}

Weighting a fit by the \emph{observed} counts biases it low; the likelihood with the
\emph{predicted} variance, or the equivalent IRLS, removes the bias. This is why our fit
iterates with $q\,\max(f,q)$. It is textbook, but the HH reweighting literature does not say
which weights its least-squares fits use.

% ===========================================================================
\section{Synthesis}
\label{sec:synth}

\begin{table}[h]
\centering\small
\begin{tabularx}{\linewidth}{@{}>{\raggedright\arraybackslash}p{4.2cm}
  >{\raggedright\arraybackslash}X >{\raggedright\arraybackslash}p{2.6cm}@{}}
\toprule
Ingredient & Closest prior work & Status \\
\midrule
23-term NLO HEFT form & Buchalla et al.\ \cite{Buchalla2018} & standard \\
Per-bin coefficients from many runs (free fit) & Carvalho \cite{Carvalho2017}, Alasfar
  \cite{Alasfar2023}, Belyaev \cite{Belyaev2025} & standard \\
Smooth functions of $x$ times known polynomial in couplings
  & Chen et al.\ \cite{Chen2021}; VCM \cite{HastieTibshirani1993} & known, not for HH \\
Spline density fitted to bin integrals of counts & PCLM \cite{Eilers2007,Rizzi2015};
  histospline \cite{Boneva1971} & known, not in HEP EFT \\
Smoothing parameters by evidence & MacKay \cite{MacKay1992}, Wood \cite{Wood2011}
  & standard in statistics \\
Exact Poisson variances from integer counts; IRLS & Baker--Cousins \cite{BakerCousins1984}
  & standard idea; count recovery is ours \\
Threshold factor $\phi(m)$ & phase-space $\beta$ factor & physics; trivial \\
Model error $\propto$ cancellation & Cadamuro et al.\ \cite{Cadamuro2025} (qualitative)
  & \textbf{not found} \\
Within-bin validation (merged-bin test) & PCLM ungrouping & \textbf{not found in HEP} \\
\textbf{All of the above, for $d\sigma/dm_{hh}$} & --- & \textbf{not found} \\
\bottomrule
\end{tabularx}
\caption{Each ingredient of the smooth model and its closest precedent.}
\label{tab:synth}
\end{table}

\begin{keyidea}[What is new]
Our model is a \textbf{varying-coefficient penalised composite link model}: the
varying-coefficient structure supplies $\sum_iC_i(\theta)a_i(m)$, the composite link supplies
the bin integrals. We found neither the combination in statistics nor any smooth
parametrisation of the HEFT coefficients in physics. Specific to this work:
\begin{enumerate}
  \item continuous $a_i(m_{hh})$ for the 23 NLO HEFT terms, with the threshold factor,
    fitted jointly to $\sim$4,000 POWHEG spectra;
  \item a likelihood on the \emph{integral} of the density over each bin, with the exact
    Poisson variances recovered from the integer event counts;
  \item 23 smoothness strengths from the evidence and an exact Gaussian posterior for every
    $a_i(m)$ and every prediction;
  \item a calibrated model error that scales with the cancellation $\sum_i|C_ia_i|$;
  \item a test of within-bin accuracy that a per-bin table cannot pass by construction.
\end{enumerate}
\end{keyidea}

\begin{watch}[What a referee might say, and our answer]
\begin{description}[leftmargin=0pt,style=nextline]
  \item[``This is just the WG per-bin fit with smoothing.''] Yes, and that is the point:
    the per-bin fit is the floor for any 23-term model, and we reach it with half the parameters
    (a third of the effective ones), continuous predictions, and a demonstrated gain inside bins and in the tail.
  \item[``Why not a network?''] We tried; an ensemble of HEFTNets reaches the same floor,
    but its uncertainty requires an ensemble or a Laplace mixture, and it is unreliable
    between bin centres (spread up to 20--90 times larger below 400~GeV).
  \item[``The discrepancy term is ad hoc.''] It has three parameters, is fitted on training
    spectra and calibrated on held-out ones (pull s.d.\ 1.00 in every cancellation fifth),
    and its form is motivated by the destructive-interference effect seen independently in
    ref.~\cite{Cadamuro2025}.
\end{description}
\end{watch}

% ===========================================================================
\section{Recommendations and actions}
\label{sec:actions}

\paragraph{For the paper.}
\begin{enumerate}
  \item In the introduction, cite \cite{Buchalla2018,Heinrich2020,Alasfar2023} as the
    per-bin baseline and say that the free $23\times80$ fit reproduces that approach with a
    much larger sample.
  \item Name the model: ``a varying-coefficient model \cite{HastieTibshirani1993} estimated
    as a penalised composite link model \cite{Eilers2007,Rizzi2015}''.
  \item Cite \cite{Chen2021,Brehmer2018} when motivating HEFTNet: polynomial in the
    couplings, learned functions of the kinematics.
  \item Cite \cite{Cadamuro2025} and \cite{Heinrich2020} (grid, Sudakov) when introducing the
    cancellation-dependent discrepancy.
  \item Offer the smooth $a_i(m)$ and their covariance as a drop-in replacement for the
    per-bin $d\textbf{A}$ tables of \cite{Alasfar2023}: same \cref{eq:alasfar}, evaluated at
    any $m_{hh}$ or integrated over any binning.
\end{enumerate}

\paragraph{Concrete follow-ups.}
\begin{enumerate}
  \item Download the 13~TeV ancillary CSV of \cite{Buchalla2018} (and the WG $d\textbf{A}$
    files if public) and overlay their $\tilde A_i$ on our $a_i(m)$ with bands, after
    matching conventions. Agreement validates both; differences in the tail show the gain.
  \item Ask whether the WG fit weights are observed or predicted variances.
  \item INSPIRE ``cited by'' search and a note to the HH conveners (see the box in section 2).
  \item Merge the entries of \file{docs/smooth/literature_refs.bib} into
    \file{docs/paper/appendix_refs.bib}.
\end{enumerate}

\begin{thebibliography}{99}\small
\bibitem{Buchalla2018} G.~Buchalla, M.~Capozi, A.~Celis, G.~Heinrich, L.~Scyboz,
  \emph{Higgs boson pair production in non-linear Effective Field Theory with full
  $m_t$-dependence at NLO QCD}, JHEP 09 (2018) 057, arXiv:1806.05162.
\bibitem{Heinrich2020} G.~Heinrich, S.\,P.~Jones, M.~Kerner, L.~Scyboz,
  \emph{A non-linear EFT description of $gg\to HH$ at NLO interfaced to POWHEG},
  JHEP 10 (2020) 021, arXiv:2006.16877.
\bibitem{Capozi2020} M.~Capozi, G.~Heinrich, \emph{Exploring anomalous couplings in Higgs
  boson pair production through shape analysis}, JHEP 03 (2020) 091, arXiv:1908.08923.
\bibitem{Alasfar2023} L.~Alasfar et al., \emph{Effective Field Theory descriptions of Higgs
  boson pair production}, LHC Higgs WG note, arXiv:2304.01968.
\bibitem{Carvalho2016} A.~Carvalho, M.~Dall'Osso, T.~Dorigo, F.~Goertz, C.\,A.~Gottardo,
  M.~Tosi, \emph{Higgs pair production: choosing benchmarks with cluster analysis},
  JHEP 04 (2016) 126, arXiv:1507.02245.
\bibitem{Carvalho2016b} A.~Carvalho et al., \emph{Analytical parametrization and shape
  classification of anomalous HH production in the EFT approach}, LHCHXSWG-2016-001,
  arXiv:1608.06578.
\bibitem{Carvalho2017} A.~Carvalho, F.~Goertz, K.~Mimasu, M.~Gouzevitch, A.~Aggarwal,
  \emph{On the reinterpretation of non-resonant searches for Higgs boson pairs},
  arXiv:1710.08261.
\bibitem{Cadamuro2025} L.~Cadamuro, T.~Ingebretsen Carlson, J.~Sj\"olin,
  \emph{Di-Higgs and Effective Field Theory: Signal Reweighting Beyond $m_{hh}$},
  arXiv:2502.20976.
\bibitem{ATLmorph} ATLAS Collaboration, \emph{A morphing technique for signal modelling in a
  multidimensional space of coupling parameters}, ATL-PHYS-PUB-2015-047.
\bibitem{RooMorph} ROOT, \code{RooLagrangianMorphFunc} class reference,
  \url{https://root.cern.ch/doc/v632/classRooLagrangianMorphFunc.html}.
\bibitem{Belyaev2025} N.~Belyaev, R.~Konoplich, K.~Prokofiev, \emph{Signal Modelling with
  Overdetermined Morphing Technique}, arXiv:2512.12473.
\bibitem{Mohrman2023} K.~Mohrman, \emph{EFT-aware histograms}, LPC EFT Workshop, Fermilab,
  6 Sept.\ 2023, \url{https://indico.fnal.gov/event/60025/}.
\bibitem{Chen2021} S.~Chen, A.~Glioti, G.~Panico, A.~Wulzer, \emph{Parametrized classifiers
  for optimal EFT sensitivity}, JHEP 05 (2021) 247, arXiv:2007.10356.
\bibitem{Brehmer2018} J.~Brehmer, K.~Cranmer, G.~Louppe, J.~Pavez, \emph{Constraining
  Effective Field Theories with Machine Learning}, Phys.\ Rev.\ Lett.\ 121 (2018) 111801,
  arXiv:1805.00013.
\bibitem{Mastandrea2024} R.~Mastandrea, B.~Nachman, T.~Plehn, \emph{Constraining the Higgs
  Potential with Neural Simulation-based Inference for Di-Higgs Production},
  arXiv:2405.15847.
\bibitem{Bahl2026} H.~Bahl, J.~Braun, G.~Heinrich, T.~Plehn, R.~Revelli, \emph{How to Trust
  Learned Loop Amplitudes}, arXiv:2601.00950.
\bibitem{HastieTibshirani1993} T.~Hastie, R.~Tibshirani, \emph{Varying-coefficient models},
  J.\ R.\ Stat.\ Soc.\ B 55 (1993) 757.
\bibitem{EilersMarx1996} P.\,H.\,C.~Eilers, B.\,D.~Marx, \emph{Flexible smoothing with
  B-splines and penalties}, Statist.\ Sci.\ 11 (1996) 89.
\bibitem{EilersMarx2021} P.\,H.\,C.~Eilers, B.\,D.~Marx, \emph{Practical Smoothing: The Joys
  of P-splines}, Cambridge University Press, 2021.
\bibitem{Eilers2007} P.\,H.\,C.~Eilers, \emph{Ill-posed problems with counts, the composite
  link model and penalized likelihood}, Statistical Modelling 7 (2007) 239.
\bibitem{Rizzi2015} S.~Rizzi, J.~Gampe, P.\,H.\,C.~Eilers, \emph{Efficient estimation of
  smooth distributions from coarsely grouped data}, Am.\ J.\ Epidemiol.\ 182 (2015) 138.
\bibitem{Boneva1971} L.\,I.~Boneva, D.~Kendall, I.~Stefanov, \emph{Spline transformations:
  three new diagnostic aids for the statistical data-analyst}, J.\ R.\ Stat.\ Soc.\ B 33
  (1971) 1.
\bibitem{MacKay1992} D.\,J.\,C.~MacKay, \emph{Bayesian interpolation}, Neural Comput.\ 4
  (1992) 415.
\bibitem{Wood2011} S.\,N.~Wood, \emph{Fast stable restricted maximum likelihood and marginal
  likelihood estimation of semiparametric generalized linear models}, J.\ R.\ Stat.\ Soc.\ B
  73 (2011) 3.
\bibitem{Whitehorn2013} N.~Whitehorn, J.~van Santen, S.~Lafebre, \emph{Penalized splines for
  smooth representation of high-dimensional Monte Carlo datasets}, Comput.\ Phys.\ Commun.\
  184 (2013) 2214, arXiv:1301.2184.
\bibitem{BakerCousins1984} S.~Baker, R.\,D.~Cousins, \emph{Clarification of the use of
  chi-square and likelihood functions in fits to histograms}, Nucl.\ Instrum.\ Meth.\ 221
  (1984) 437.
\end{thebibliography}

\end{document}
